\mychapter{Metodologia}{cap:metod}\lhead{METODOLOGIA}

% \section*{Resumo executivo}
% Este capítulo descreve, de forma operacional e reprodutível, a metodologia para avaliar soluções de MEIO (\emph{Multi-Echelon Inventory Optimization}) em rede geral (fábricas--centros--lojas), combinando previsão de demanda, otimização (ILP/MILP e metaheurísticas) e políticas aprendidas (DRL/MARL) em simulação sob \emph{rolling horizon}. \cite{Shapiro2001ModelingSupplyChain,AhujaMagnantiOrlin1993NetworkFlows,BirgeLouveaux2011StochasticProgramming,Law2015SimulationModeling,Boute2022DRLInventoryRoadmap,Geevers2024MEIODRL,Ziegner2025IMARL}\footnote{Capítulo elaborado em 15/03/2026 (America/Maceio).}

\section{Objetivos experimentais, hipóteses e suposições}
\subsection*{Metas rastreáveis aos OE1--OE7}
O desenho experimental valida: (OE1) consistência do modelo base em rede; (OE2) geração de instâncias controladas e uso de dados de grande escala (conforme apresentação do autor); (OE3) \emph{benchmark} ILP/MILP; (OE4) escalabilidade por GA/GRASP/SA e, quando multiobjetivo, NSGA-II; (OE5--OE6) extensão multi-período/estocástica em \emph{rolling horizon}; (OE7) políticas DRL/MARL (PPO/DQN/IMARL) em ambiente simulado.

\subsection*{Hipóteses}
(H1) ILP/MILP fornece referência com \emph{gap} controlado em instâncias pequenas/médias, mas escala de forma limitada com \(|A|\times|I|\times|T|\). \cite{Shapiro2001ModelingSupplyChain}
(H2) Em instâncias grandes, metaheurísticas atingem soluções competitivas com menor \emph{runtime} e robustez sob múltiplas sementes. \cite{Talbi2009Metaheuristics}
(H3) Em multi-período estocástico, DRL/MARL supera abordagens estáticas ao aprender políticas adaptativas, desde que o ambiente/recompensa estejam bem definidos. \cite{Boute2022DRLInventoryRoadmap,Geevers2024MEIODRL,Ziegner2025IMARL}
(H4) A previsão de demanda explica variação relevante em custo e nível de serviço; métricas alinhadas à decisão são mais informativas que erro puro. \cite{Nasseri2023RetailDemandPrediction,TadayonradNdiaye2023KPIDemandForecasting}

\subsection*{Suposições explícitas (itens não especificados)}
Horizon \(H\) (propor \(H\in\{14,28,56\}\) dias ou \(H\in\{8,12\}\) semanas), política de ruptura (backorder \emph{ou} lost sales), lead time \(\ell_{a,b}\) (determinístico \emph{ou} discreto em \(\{0,1,2,3\}\)) e ruído de demanda (coeficiente \(cv\in\{0,0.2,0.5\}\)). \cite{Axsater2006InventoryControl,BirgeLouveaux2011StochasticProgramming}

\section{Dados, pré-processamento e geração de instâncias}
\subsection{Dados e esquema mínimo}
Conforme apresentação do autor, o caso motivador possui escala aproximada de \(48\)M registros, \(12\)k produtos únicos e \(3\)M entidades de demanda, implicando necessidade de agregação e/ou amostragem em experimentos controlados. No pipeline Kedro anexado, camadas \texttt{01\_raw}--\texttt{08\_reporting} organizam limpeza, enriquecimento e previsão; a tabela de vendas enriquecida inclui ao menos \texttt{date}, \texttt{cycle}, \texttt{qty}, receitas, canais, chaves \texttt{rev\_id}/\texttt{prod\_id} e atributos de revendedor/produto; o módulo de previsão compara \emph{naive}, Holt-Winters, SARIMAX, XGBoost e Prophet sob validação \emph{rolling-origin} e métricas MAE/RMSE/MAPE (conforme implementação). \cite{Nasseri2023RetailDemandPrediction}

\subsection{Pré-processamento e particionamento temporal}
As transformações incluem padronização de chaves, remoção de PII, deduplicação e agregação por ciclo (parâmetros em YAML; semente global \texttt{random\_seed}=123 conforme configuração). O particionamento treino/val/test é temporal: treino até \(t_{\mathrm{train}}\), validação em \(t_{\mathrm{train}}{+}1{:}t_{\mathrm{val}}\) e teste nos últimos \(H\) ciclos (não especificado). O mesmo princípio é aplicado à avaliação de políticas em \emph{rolling horizon} para evitar vazamento de informação futura.

\subsection{Geração de instâncias}
Instâncias controladas são definidas por \(|M|,|S|,|F|,|I|\), horizonte \(H\) e ruído \(cv\). Protocolo: amostrar subconjuntos de nós/itens; gerar distâncias \(d(a,b)\) (não especificado; p.ex.\ coordenadas 2D) e filtrar arcos por viabilidade \(d(a,b)\le T_{a,b}\); definir \(C_{a,b}\) e \(Q_f\) proporcionais à demanda média; gerar demanda base e cenários \(D_{m,i,t}\) via (i) bootstrap do histórico ou (ii) \(D=\mu(1+\epsilon)\) com \(\epsilon\sim \mathcal{N}(0,cv)\) truncado. \cite{AhujaMagnantiOrlin1993NetworkFlows,Shapiro2001ModelingSupplyChain}

\section{Métodos avaliados e parametrização}
\subsection{ILP/MILP (benchmark) e solver}
O baseline resolve o problema de um período com \(x_{a,b,i}\) e \(y_{f,i}\) e objetivo com penalidade temporal \(\lambda_t\) (grade não especificada; sugerir \(\{0,0.1,1,10\}\) após normalização). \cite{Shapiro2001ModelingSupplyChain} Configurações padrão: \texttt{TimeLimit}=300\,s, \texttt{MIPGap}=1\%, \texttt{Threads}=1; reportam-se \emph{incumbent}, \emph{best bound} e \emph{gap}.

\subsection{Metaheurísticas (GA/GRASP/SA/NSGA-II)}
GA (codificação ``1 gene por item'') minimiza custo de transporte + produção adicional; parâmetros iniciais \(ngen{=}50\), \(pop{=}80\), \(cxpb{=}0.7\), \(mutpb{=}0.2\), torneio \(=3\) (conforme implementação e apresentação do autor). GRASP usa construção randomizada com \(\alpha\) (propor \(\alpha\in\{0.1,0.3,0.5\}\)) e busca local; SA usa resfriamento geométrico \(T_{k+1}=\alpha T_k\) (propor \(\alpha{=}0.95\)). \cite{Talbi2009Metaheuristics}
NSGA-II é usado quando custo e tempo são objetivos separados (alternativa ao ajuste de \(\lambda_t\)), produzindo aproximação empírica de fronteiras de Pareto.

\subsection{Simulação (Monte Carlo / eventos discretos)}
A avaliação multi-período é feita por simulação: para cada replicação \(r=1,\dots,R\) (\(R\) não especificado; sugerir \(R=30\)), amostram-se demandas e (opcionalmente) lead times; aplica-se decisão por período via MILP/heurística ou política RL; acumula-se custo e nível de serviço. Para reduzir variância comparativa, usa-se \emph{common random numbers} (mesmas sementes/cenários por método). \cite{Law2015SimulationModeling,BirgeLouveaux2011StochasticProgramming}

\subsection{DRL/MARL (PPO/DQN/IMARL)}
O ambiente é modelado como MDP; ações representam reposição/transferências (vetor contínuo ou discretizado). PPO é o baseline de política em rede geral (treino com Stable-Baselines3); DQN é avaliado apenas quando o espaço de ação é discretizado; IMARL é considerado como extensão multiagente voltada à escalabilidade em MEIO realista. \cite{Schulman2017PPO,Oroojlooyjadid2022BeerGameDQN,Geevers2024MEIODRL,Ziegner2025IMARL}

\section{Ambiente de simulação e API tipo Gym}
\subsection{Estado, ação, transição e recompensa}
O estado \(s_t\) inclui estoques \(I_{n,i,t}\), pipeline (pedidos em trânsito), demanda recente/prevista e capacidades residuais; a ação \(a_t\) define \(x_{a,b,i,t}\) e \(y_{f,i,t}\) (diretamente ou via parametrização). A transição aplica lead time \(\ell\), realiza demanda \(D_{m,i,t}\) e atualiza balanços; a recompensa é \(-\Delta\)custo, incluindo transporte, produção, holding, ruptura e termo temporal ponderado por \(\lambda_t\). \cite{Axsater2006InventoryControl,Boute2022DRLInventoryRoadmap}

\subsection{Rolling horizon}
Para cada época \(t\): consolidar \([t-w,t]\) (\(w\) não especificado), prever \(\hat{D}_{t+1:t+H}\), decidir para o próximo período (MILP/heurística ou RL), executar e registrar novos dados; repetir até \(H\). \cite{Law2015SimulationModeling}

\section{Protocolo experimental, métricas e testes}
\subsection{Sementes e splits}
Todas as execuções registram sementes (Python/NumPy/ambiente) e, quando necessário, fixam paralelismo. Previsão e avaliação são temporais; DRL treina em instâncias de treino e reporta desempenho em instâncias \emph{hold-out}. Para comparações estocásticas, utiliza-se CRN (mesmas amostras por método) e múltiplas seeds.

\subsection{Métricas}
Operacionais: custo total e decomposição, fill rate, backorders/lost sales, estoque médio, violações; computacionais: runtime, memória (quando disponível), \emph{gap} (MILP), retorno médio e estabilidade por sementes (RL). \cite{Axsater2006InventoryControl,Law2015SimulationModeling}

\subsection{Inferência}
Comparações pareadas usam Wilcoxon (ou t-teste se adequado), com \(\alpha=0.05\); intervalos de confiança por bootstrap são opcionais. \cite{Law2015SimulationModeling}

\section{Reprodutibilidade e engenharia experimental}
O pipeline é orquestrado em Kedro (pipelines \texttt{data\_clean}, \texttt{data\_enrich}, \texttt{forecasting}, \texttt{reporting}) com parâmetros em YAML e rastreamento de métricas/artefatos via MLflow (conforme código anexado). A reprodutibilidade é garantida por: versionamento Git (tags por experimento), hashes de datasets/configs, logs estruturados (JSON) e containerização via Docker (dependências fixadas).

\subsection*{Diagrama: fluxo de dados e de experimentos (Mermaid)}
\begin{verbatim}
flowchart TD
  R[01_raw: vendas/ruptura] --> C[data_clean]
  C --> E[data_enrich: fatos + dimensões]
  E --> F[forecasting: MAE/RMSE/MAPE + \hat{D}]
  F --> O[otimização: MILP + metaheurísticas]
  O --> S[simulação/RL: rolling horizon]
  S --> P[08_reporting: métricas + gráficos]
\end{verbatim}

\begin{verbatim}
flowchart LR
  A[Instâncias × seeds] --> B[MILP]
  A --> C[GA/GRASP/SA/NSGA-II]
  A --> D[Simulação (R rep.)]
  D --> E[RL treino]
  B --> Z[Relatório]
  C --> Z
  E --> Z
\end{verbatim}

\section{Tabelas e matriz experimental}
\begin{table}[htbp]\centering\small
\begin{tabular}{lcccc}
\hline
Método & Escala & Interpret. & Custo comp. & Restrições \\
\hline
MILP & M--B & Alta & Alto & Exato \\
GA/GRASP/SA & Alta & Média & Médio & Penal./reparo \\
Simulação & Alta & Alta & Médio--alto & Avalia \\
PPO/DQN/IMARL & M & Baixa & Alto & Indireto \\
\hline
\end{tabular}
\caption{Comparação resumida de famílias (MEIO).}\label{tab:met_comp}
\end{table}

\begin{table}[htbp]\centering\small
\begin{tabular}{lccc}
\hline
Classe & \((|M|,|S|,|F|,|I|)\) & Métodos & Seeds/Rep. \\
\hline
S & (20,2,1,50) & MILP+meta & 10 \\
M & (50,5,2,200) & MILP+meta & 10 \\
L & (100,10,5,500) & meta+RL & 10 / \(R{=}30\) \\
\hline
\end{tabular}
\caption{Matriz recomendada; \(H\) e \(cv\) varrem cenários.}\label{tab:met_matrix}
\end{table}

\section*{Snippets de implementação (documentação)}
\subsection*{OR-Tools/Gurobi (esqueleto MIP)}
\begin{verbatim}
# declarar solver/model; criar x[a,b,i], y[f,i]
# add: atendimento, capacidade, produção; objetivo com lambda_t
# set TimeLimit/MIPGap/Threads; resolver; salvar incumbent/bound/gap
\end{verbatim}

\subsection*{GA/GRASP (pseudocódigo)}
\begin{verbatim}
GA: população -> avaliar -> seleção(torneio) -> crossover(2pt) -> mutação -> elitismo.
GRASP: construir (RCL com α) -> busca local -> atualizar melhor; repetir K iterações.
\end{verbatim}

\subsection*{Gym-like env (esqueleto)}
\begin{verbatim}
reset(seed): t=0; inicializa estoques/pipeline; retorna obs.
step(a): aplica decisão; amostra demanda/lead time; atualiza estado e custo;
return obs, reward=-custo_inc, terminated=(t>=H).
\end{verbatim}

\subsection*{PPO (loop)}
\begin{verbatim}
model = PPO("MlpPolicy", env, n_steps=2048, batch_size=128, gamma=0.99, lr=3e-4)
model.learn(total_timesteps=2_000_000); avaliar em seeds e instâncias hold-out.
\end{verbatim}

\subsection*{Kedro node (exemplo)}
\begin{verbatim}
def prepare_series(vendas_enrich):
    return vendas_enrich.groupby("cycle")["qty"].sum().sort_index()
\end{verbatim}

\subsection*{Dockerfile (esboço)}
\begin{verbatim}
FROM python:3.11-slim
WORKDIR /app
COPY requirements.txt ./
RUN pip install -r requirements.txt
COPY . .
CMD ["kedro","run"]
\end{verbatim}
